%===========================================
% 第11章 結論（限界と今後の課題）
% 研究結果をまとめ、理論的・実証的限界を整理し、今後の研究課題（データ精緻化、動学モデル化、他産業への応用可能性）を示す。
%===========================================


%%%%%%%%%%%%%%%%%%%%%%%%%%%%%%%%%%%%%%%%%%%%%%%%%%%%%%%%%%%%%%%%%%%%%%%%%%%%%%
% 【研究の中心命題（絶対に忘れてはならない論文の背骨）】
%
% ● JEPX のスポット価格は平常時は安定している。
% ● 2021 年 1 月の高騰は「単発のショック」にすぎない。
% ● にもかかわらず、多くの新電力が倒産した。
%
% → これは「価格」では倒産を説明できないことを意味する。
%
% 【本研究の核心】
% 倒産の本質的原因は、価格ショックそのものではなく、
% 「新電力ビジネスモデルの構造的脆弱性」にある。
%
% 【構造的脆弱性の主要因】
% 1. 先渡し契約（ヘッジ）の未整備 → スポット依存 → 逆ザヤ発生
% 2. 資本力の弱さ → 数週間の逆ザヤで資金ショート
% 3. 需給調整市場コストの増加 → 小規模ほど負担が重い
% 4. 需要予測の精度不足 → 追加調達で高値掴み
% 5. JEPX の担保制度強化 → 小規模事業者は担保を積めず退出
%
% 【重要な論点】
% ● 2021 年の高騰は「引き金（トリガー）」にすぎない。
% ● 倒産の根本原因は「構造」にある。
% ● 公開データ（価格・取引量・参加者数）だけでは倒産は説明できない。
%
% 【本研究の価値（独自性）】
% ● 公開データでは見えない倒産メカニズムを解明する。
% ● 新電力の財務構造 × 市場制度 × 価格ショックの相互作用を分析する。
% ● 「倒産して当然」ではなく、「倒産する構造だった」ことを示す。
%
% 【論文全体のストーリーライン】
% 1. 市場データ分析（価格・取引量・参加者数）
%    → 表面的には安定しており、倒産は説明できない。
%
% 2. 新電力のビジネスモデル分析
%    → 構造的脆弱性を明確化。
%
% 3. 2021 年高騰の影響分析
%    → 価格ショックは「構造的弱点を露呈させた」だけ。
%
% 4. 倒産リスクの総合モデル化
%    → 財務指標 × 市場構造 × 価格ショック。
%
% 5. 政策的含意
%    → ヘッジ義務、担保制度、調整市場設計など。
%
% 【最終結論の方向性】
% 「新電力の倒産は、単発の価格ショックではなく、
% 　ビジネスモデルの構造的脆弱性が根本原因である」
%
% この方向性に論文全体を収束させること。
%%%%%%%%%%%%%%%%%%%%%%%%%%%%%%%%%%%%%%%%%%%%%%%%%%%%%%%%%%%%%%%%%%%%%%%%%%%%%%

% ------------------------------------------------------------
% 【研究の落としどころ（結論の核心）に関するメモ】
% ------------------------------------------------------------
%
% 本研究の結論は、「新電力企業の倒産は、財務指標だけでは説明できない
% “構造的リスク”によって決まる」という点にある。
%
% これは自明ではなく、また奇をてらった結論でもない。
% むしろ、第2章の倒産事例観察、第3〜5章の理論構築、
% 第6章の相互作用分析、第7章の実証分析が自然に導く帰結である。
%
% 【1. 財務指標では説明できない倒産】
% 新電力企業の倒産は、従来の倒産モデル（Altman Z、TSR/TDB、Merton、B&S）
% が前提とする「財務比率」や「収益性」だけでは説明できない。
% F-Power やウエスト電力の事例が示すように、
% 財務諸表上は急激な悪化が見えない段階でも倒産が発生している。
%
% 【2. 倒産を決定づけたのは μ・σ・S・D の“構造要因”】
% 本研究が形式化した μ（収益性）、σ（外部ショック感応度）、
% S（支援能力）、D（倒産閾値）の四要素は、
% 新電力市場の制度構造・市場構造・組織構造を反映している。
%
% 特に σ・S・D は財務諸表には現れないが、
% 倒産を決定づけた本質的な要因である。
% これらを明示的にモデル化した点が本研究の独自性である。
%
% 【3. 従来モデルには存在しない構造的メカニズム】
% Merton/B&S は株価と負債構造を前提とするが、
% 新電力企業は株価データを持たず、制度コストが外生的に変動する。
% 本研究はこの制度的特徴を μ・σ・S・D として取り込み、
% 新電力特有の倒産急増を説明する枠組みを提供する。
%
% 【4. Z_lin は“新電力専用の倒産モデル”として意味を持つ】
% 本研究の潜在倒産距離 Z_lin は、
% 新電力企業の倒産メカニズムを構造的に捉えるための指標であり、
% 従来の財務比率ベースのモデルでは捉えられない倒産リスクを説明する。
%
% 【5. 結論の位置づけ】
% 以上より、本研究の結論は次のように整理される：
%
% 「新電力企業の倒産は、財務指標ではなく、
% 市場構造・制度構造・支援構造によって決まる。
% これらを μ・σ・S・D の四要素として形式化することで、
% 従来モデルでは説明できなかった倒産急増のメカニズムを
% 理論的にも実証的にも説明できる。」
%
% この結論は自明ではなく、奇をてらったものでもなく、
% 本研究の全章が自然に導く“唯一の落としどころ”である。
%
% ------------------------------------------------------------
\iffalse
# ⭐ **第10章が「構造的解釈」なら、第11章は  
“研究全体の総括・限界・政策含意・今後の課題” に専念する章になる。**

つまり、第11章は **第9章の結果解釈をする章ではなくなる**。  
第10章で結果の意味を十分に解釈するため、第11章は「総括」に集中できます。

---

# ◆ 第10章と第11章の役割分担（最終版）

| 章 | 役割 |
|----|------|
| **第9章** | Ordered Logit の推定結果（事実の提示） |
| **第10章** | 推定結果の構造的解釈（理論 × 実証の統合） |
| **第11章** | 研究全体の結論・限界・政策含意・今後の課題 |

この三段構造は論文として非常に美しい流れになります。

---

# ⭐ **第11章 結論（限界と今後の課題）— 改訂版（仮説回収を統合済）**

### **11.1 本研究の総括：中心命題の再確認**

- JEPX のスポット価格は平常時は安定  
- 2021 年 1 月の高騰は「単発ショック」にすぎない  
- にもかかわらず新電力が大量倒産した  
- → 価格では倒産を説明できない  
- → 倒産の本質は「構造的脆弱性」にある  

ここで **論文の背骨（中心命題）を再提示**します。

---

### **11.2 本研究が明らかにした構造的脆弱性**

- ヘッジ未整備 → スポット依存 → 逆ザヤ  
- 資本力の弱さ → 数週間の逆ザヤで資金ショート  
- 需給調整市場コスト → 小規模ほど負担が重い  
- 需要予測の精度不足 → 高値掴み  
- JEPX の担保制度強化 → 小規模事業者は退出  

ここは **第2章〜第5章の理論構築と、第10章の構造的解釈を統合する部分**です。

---

### **11.3 本研究の理論的貢献（独自性）**

- 公開データでは見えない倒産メカニズムを形式化  
- μ・σ・S・D の四要素が新電力の制度構造を反映  
- 従来モデル（Altman Z、TSR/TDB、Merton/B&S）では説明できない倒産を説明  
- 潜在倒産距離 \(Z_{\text{lin}}\) を新電力専用モデルとして提示  

ここは **あなたの研究の価値を明確に示す部分**です。

---

### **11.4 実証分析の結論：日本企業と英国企業の比較**

- 日本企業：日次データ → σ は std 一択  
- 英国企業：月次データ → σ は 3 指標でロバスト性検証  
- Ordered Logit により倒産ステージの段階性を確認  
- μ・σ・D が倒産ステージを押し上げる主要因  
- S（支援能力）は倒産回避に重要（日本企業で顕著）  
- → 国際比較により「構造的脆弱性」が共通要因であることを確認  

ここは **第9章・第10章の結果を総括する部分**です。

---

# ⭐ **11.5 提示した仮説の妥当性の検討（H1〜H4）**

以下に、先ほど作成した LaTeX テンプレートをそのまま挿入できます。

（※ここでは構成案のみ示し、本文は省略します  
→ 実際の本文は前回お渡しした LaTeX テンプレートを貼り付けてください）

---

### **11.6 政策的含意**

- ヘッジ義務の導入  
- 担保制度の再設計  
- 小規模新電力の資本規制  
- 需給調整市場コストの再評価  
- 親会社支援の透明化  
- 新電力のリスク管理指針の整備  

ここは **あなたの研究の社会的意義を示す部分**です。

---

### **11.7 本研究の限界**

- 日次データの制約（英国）  
- S（支援能力）の代理変数の限界  
- D（倒産閾値）の代理変数の単純化  
- Ordered Logit の静学モデルとしての限界  
- 企業行動の内生性を扱えない点  

ここは **審査者が必ず見るポイント**です。

---

### **11.8 今後の研究課題**

- データ精緻化（英国の日次データ取得）  
- 動学モデル化（時系列依存を組み込む）  
- 他産業への応用（ガス小売、通信など）  
- 親会社支援の定量化  
- 企業行動の内生性を扱う構造推定モデルへの拡張  

ここは **あなたの研究の発展可能性を示す部分**です。

---

### **11.9 最終結論**

> 「新電力企業の倒産は、単発の価格ショックではなく、  
> 　ビジネスモデルの構造的脆弱性が根本原因である。」


\fi

% ------------------------------------------------------------
% 11.5 提示した仮説の妥当性の検討
% ------------------------------------------------------------

\section{提示した仮説の妥当性の検討}

本研究では，第3章において新電力企業の倒産メカニズムに関する
4 つの仮説（H1～H4）を提示した。本節では，
第9章の Ordered Logit 推定結果および第10章で行った構造的解釈に基づき，
各仮説が支持されるか否かを総合的に評価する。

\subsection*{H1：収益性（$\mu$）は倒産ステージを押し上げる主要因である}

第9章の推定結果では，資産成長率（$\mu$）の係数は正であり，
統計的に有意であった。これは，収益性の低下が倒産ステージの進行を
促進することを示している。また，第10章で示したように，
新電力企業のビジネスモデルは収益性の変動に対して脆弱であり，
制度コストやスポット依存構造がその影響を増幅する。
以上より，\textbf{H1 は支持される}。

\subsection*{H2：外部ショック感応度（$\sigma$）は倒産ステージを押し上げる}

日本企業では日次データに基づく標準偏差（std）が有意であり，
英国企業でも月次データを用いた 3 種類の代理変数のうち複数が
倒産ステージに対して有意な影響を示した。
第10章で示したように，新電力企業はスポット依存構造により
価格ショックの影響を直接受けるため，$\sigma$ が倒産リスクを
押し上げる構造的理由が存在する。
以上より，\textbf{H2 は支持される}。

\subsection*{H3：支援能力（$S$）は倒産回避に寄与する}

日本企業の分析では，親会社の支援能力を表す $S$ が倒産ステージを
有意に押し下げる効果を示した。一方，英国企業では公開データの制約により
$S=0$ として扱ったが，第10章で示したように，支援能力の有無は
新電力企業の倒産回避に構造的に重要である。
以上より，\textbf{H3 は日本企業において支持され，英国企業については
データ制約により検証不能であるが，構造的には支持される}。

\subsection*{H4：倒産閾値（$D$）は倒産ステージを規定する主要因である}

負債比率（$D$）は，日本企業・英国企業の双方で倒産ステージに対して
正の有意な影響を示した。これは，第5章で理論的に導出した
倒産閾値の概念と整合的であり，負債の大きさが倒産境界を押し上げる
構造的メカニズムが実証的にも確認されたことを意味する。
以上より，\textbf{H4 は支持される}。

\subsection*{総合評価}

以上の仮説回収により，本研究の中心命題である

\begin{quote}
「新電力企業の倒産は，単発の価格ショックではなく，
ビジネスモデルの構造的脆弱性によって決まる」
\end{quote}

という結論が，理論的にも実証的にも支持されることが確認された。

% ============================================================
% 11.7 本研究の限界
% ============================================================

\section{本研究の限界}

本研究は，新電力企業の倒産プロセスを潜在変数モデルと Ordered Logit によって
形式化し，日本企業と英国企業の比較を通じて構造的脆弱性を明らかにした。
しかし，以下の点において本研究には限界が存在する。

\subsubsection*{（1）データ頻度の制約}

日本企業については，JEPX の日次スポット価格データ（Day-Ahead）を用いることができた一方，
英国企業については，英国卸電力市場のスポット価格（Day-Ahead Baseload）の
月次データしか利用できず，ショック感応度 $\sigma$ の推定精度に差が生じた。
英国企業の倒産プロセスをより精緻に分析するためには，
日次のスポット価格データや需給調整市場（Imbalance market）の高頻度データを取得することが望まれる。

さらに，収益性 $\mu$ についてもデータ制約が存在する。
日本企業では公開財務諸表の制約から P/L データが利用できず，
B/S の Equity/TotalAssets を収益性の代理変数として用いる必要があった。
一方，英国企業では P/L が公開されており，
本来の利益率（Operating profit や Profit before tax）を
収益性 $\mu$ として直接利用できるため，
今後はより精緻な収益性指標を導入することで，
潜在変数モデルの推定精度を向上させる余地が大きい。

\subsubsection*{（2）支援能力 $S$ の代理変数の限界}

支援能力 $S$ は倒産回避に重要な役割を果たすが，
本研究では公開財務データを用いた代理変数によって推定した。
親会社支援の実態や資本注入のタイミングなど，
より詳細な企業内部データを用いることで，
$S$ の効果をより正確に測定できる可能性がある。

\subsubsection*{（3）倒産閾値 $D$ の単純化}

倒産閾値 $D$ は企業の財務構造を反映する重要な要素であるが，
本研究では負債依存度を中心とした単純化された指標を用いた。
実際には，資金繰り，信用供与，担保制度，取引先との関係など，
複数の要因が複合的に倒産閾値を規定している可能性がある。

\subsubsection*{（4）Ordered Logit の静学モデルとしての限界}

Ordered Logit モデルは，倒産ステージの段階性を捉える上で有効である一方，
時間依存性を扱うことができない静学モデルである。
倒産プロセスは本来，企業行動や市場環境の変化が累積する動学的現象であり，
時系列依存を組み込んだモデル（動学 Ordered Logit や状態空間モデル，
マルコフ転移モデルなど）がより適合的である可能性がある\footnote{
動学 Ordered Logit は，潜在変数やステージが前期の状態に依存して変化する
「時系列依存」を組み込んだ拡張モデルであり，倒産ステージの推移を
マルコフ過程として扱うことができる。

状態空間モデル（State-space model）は，観測できない潜在状態
（財務健全性，ショック蓄積など）を時系列で推定する枠組みで，
カルマンフィルタ等を用いて企業の内部状態の推移を動学的に推定できる。

マルコフ転移モデル（Markov switching model）は，企業が複数の
「状態（健全・脆弱・危機など）」を確率的に切り替える構造を持ち，
外生ショックや制度変更によって状態が転移する現象を捉えることができる。

これらの動学モデルは，静学 Ordered Logit では扱えない
「時間依存性」「ショックの累積」「状態の遷移」を明示的にモデル化できるため，
倒産プロセスの理解を大きく深める可能性がある。
}.


\subsubsection*{（5）企業行動の内生性を扱えない点}

本研究のモデルは，企業行動（ヘッジ戦略，需給予測，資金繰り管理など）が
外生的に与えられることを前提としている\footnote{
企業行動の内生性とは，企業の意思決定（ヘッジ比率，需給予測精度，資金繰り方針など）が，
市場環境（価格変動，制度変更）や財務状態（資本力，負債構造）と相互作用しながら
内生的に変化することを指す。例えば，価格変動が激しくなれば企業はヘッジ比率を高める可能性があり，
逆に資金繰りが逼迫すればヘッジを縮小するなど，企業行動は外生的に固定されているわけではない。
このため，企業行動を外生的に扱うモデルでは，倒産プロセスの一部が過度に単純化される可能性がある。
}。
しかし，実際には企業行動は市場環境や財務状態と相互作用しながら変化する。
企業行動の内生性を扱う構造推定モデルを導入することで，
倒産プロセスの理解がさらに深まる可能性がある。

\subsubsection*{総括}

以上の限界は，本研究の射程を明確化するものであり，
本研究の中心命題を損なうものではない。
むしろ，構造的脆弱性という枠組みをより精緻化し，
新電力企業の倒産メカニズムを多面的に理解するための
今後の研究課題を示すものである。


% ============================================================
% 11.8 今後の研究課題
% ============================================================

\section{今後の研究課題}

本研究は，新電力企業の倒産プロセスを潜在変数モデルと Ordered Logit によって
形式化し，日本企業と英国企業の比較を通じて構造的脆弱性を明らかにした。
しかし，11.7節 で述べた限界は，本研究の射程を明確化すると同時に，
今後の研究の発展可能性を示すものである。
以下では，これらの限界を踏まえた今後の研究課題を整理する。

\subsubsection*{（1）収益性 $\mu$ の精緻化と財務データの拡張}

本研究では，日本企業・英国企業ともに公開財務データの制約から，
収益性 $\mu$ を B/S ベースの代理変数（Equity/TotalAssets）によって推定した。
しかし，本来の収益性は P/L に基づく利益率（Operating profit margin，
Profit before tax など）によって測定されるべきである。

日本企業では，公開財務諸表の制約により P/L データが利用できず，
代理変数による推定に留まった。一方，英国企業では Companies House において
P/L が公開されており，今後は利益率を直接 $\mu$ として導入することで，
潜在変数モデルの推定精度を大幅に向上させる余地がある。
日本企業についても，財務データの拡充や補完的データソースの活用により，
より精緻な収益性指標を構築することが今後の重要な研究課題となる。

\subsubsection*{（2）ボラティリティ指標 $\sigma$ の高頻度化と精緻化（日本・英国共通）}

日本企業については JEPX の日次スポット価格データを利用できた一方，
英国企業については英国卸電力市場のスポット価格（Day-Ahead Baseload）の
月次データしか利用できず，ショック感応度 $\sigma$ の推定精度に差が生じた。
日本・英国の双方において，日次スポット価格データや需給調整市場
（Imbalance market）の高頻度データを導入することで，
外部ショック感応度の推定精度を高め，倒産プロセスの動学的特徴を
より詳細に分析できる可能性がある。

\subsubsection*{（3）支援能力 $S$ の定量化と企業内部データの活用}

支援能力 $S$ は倒産回避に重要な役割を果たすが，
本研究では公開財務データを用いた代理変数によって推定した。
親会社支援の実態，資本注入のタイミング，信用供与の履歴など，
企業内部データを活用することで，$S$ の効果をより精緻に測定できる。
特に，日本企業では $S$ の影響が大きいため，
支援能力の定量化は今後の重要な研究課題である。

\subsubsection*{（4）倒産閾値 $D$ の多面的構造の解明}

倒産閾値 $D$ は負債依存度だけでなく，
資金繰り，担保制度，信用供与，取引先との関係など，
複数の要因が複合的に規定している可能性がある。
これらの要因を統合した多面的な倒産閾値モデルを構築することで，
潜在変数 $y^*$ の精度向上が期待できる。

\subsubsection*{（5）動学モデルへの拡張}

Ordered Logit は倒産ステージの段階性を捉える上で有効である一方，
時間依存性を扱うことができない静学モデルである。
倒産プロセスは本来，企業行動や市場環境の変化が累積する動学的現象であるため，
状態空間モデル，動学 Ordered Logit，マルコフ転移モデルなど，
時間依存性を組み込んだ動学モデルへの拡張が求められる。

\subsubsection*{（6）企業行動の内生性を扱う構造推定モデルの導入}

本研究のモデルは，企業行動（ヘッジ戦略，需給予測，資金繰り管理など）が
外生的に与えられることを前提としている.
しかし，実際には企業行動は市場環境や財務状態と相互作用しながら変化する。
企業行動の内生性を扱う構造推定モデルを導入することで，
倒産プロセスの理解がさらに深まり，政策評価にも応用可能となる。

\subsubsection*{（7）他産業への応用可能性の検討}

本研究で構築した潜在変数モデルと閾値構造は，
新電力企業に特化した枠組みであるが，
ガス小売，通信，地域金融など，
外生ショックと制度構造の影響を強く受ける産業にも応用可能である。
産業横断的な比較分析を行うことで，
構造的脆弱性という概念の一般性を検証できる。


\subsubsection*{総括}

以上の研究課題は，本研究の限界を踏まえつつ，
新電力企業の倒産メカニズムをより精緻に理解するための方向性を示すものである。
構造的脆弱性という枠組みは，新電力市場に限らず，
外生ショックに対して脆弱な産業を分析するための有力な視座となり得る。

% ============================================================
% 11.9 最終結論
% ============================================================

\section{最終結論}

新電力企業の倒産は，価格ショックの大小ではなく，
本研究で定式化した財務要因（収益性 $\mu$，倒産閾値 $D$）と
市場要因（ショック感応度 $\sigma$），支援能力 $S$ の構造的組み合わせによって
形成されるリスク構造の中で発生する現象として理解できる。
収益性の低下，外部ショックへの高い感応度，負債構造に起因する閾値の上昇，
支援能力の不足といった要因が重層的に作用することで，
潜在変数 $y^*$ は倒産 cutpoint 近傍へと押し上げられ，
企業は外生ショックに対して極めて脆弱な状態に置かれる。

日本企業と英国企業の比較分析（第9章〜第10章）では，
制度環境が異なるにもかかわらず，倒産リスクが
財務構造と市場制度の相互作用によって規定されるという点で
共通のメカニズムを持つことが確認された。
すなわち，新電力企業の倒産は国や市場を問わず，
同様のリスク構造のもとで段階的に進行する現象である。

さらに，本研究は，新電力企業の倒産を二値的イベントとしてではなく，
潜在倒産距離 \(Z_{\text{lin}}\) と Ordered Logit の閾値構造を用いて
ステージの推移として捉える理論的枠組みを提示した。
これにより，倒産プロセスを財務構造と市場制度の相互作用として理解する
新たな分析視座を提供した点に，本研究の独自性がある。

以上の結果を総合すると，本研究の最終結論は次のように整理される。

\begin{quote}
新電力企業の倒産は，価格ショックそのものではなく，
ビジネスモデルに内在する構造的脆弱性によって規定される。
倒産は突然発生するのではなく，制度構造と財務構造の複合的作用によって
段階的に進行する非線形現象である。
\end{quote}

本研究は，新電力企業の倒産メカニズムを制度構造と財務構造の観点から
統合的に理解するための基礎的枠組みを提示した。
この枠組みは，新電力市場の持続可能性を高める政策設計に資するだけでなく，
外生ショックに対して脆弱な他産業の分析にも応用可能である。
